\section*{Capítulo \ref{ch:g3:measuring-mass} --- Medir la masa}

\begin{solution}{exo:g3:measuring-mass:1}
De cuánta sustancia está hecha una cosa. \omterm{def:g2:measuring-length:measure}{Unidades}: el \omterm{def:g3:measuring-mass:units}{gramo} y el
\omterm{def:g3:measuring-mass:units}{kilogramo}, con $\qty{1}{kg} = \qty{1000}{g}$.
\end{solution}

\begin{solution}{exo:g3:measuring-mass:2}
Clip: \qty{1}{g}. Manzana: \qty{150}{g}. Cartón de leche:
\qty{1}{kg}. Gato: \qty{4}{kg}.
\end{solution}

\begin{solution}{exo:g3:measuring-mass:3}
$500 + 200 + 200 = 900$: la \omterm{def:g3:measuring-mass:mass}{masa} de la calabaza es \qty{900}{g}.
\end{solution}

\begin{solution}{exo:g3:measuring-mass:4}
Si la \omterm{def:g3:levers-and-scales:scale}{balanza} vacía ya se inclina, todas las pesadas mienten a favor
de un lado --- como el brazo largo del comerciante tramposo. Quedar
horizontal en vacío es la prueba de que se empieza con justicia.
\end{solution}

\begin{solution}{exo:g3:measuring-mass:5}
El balón de playa es grande pero es casi todo \omterm{def:g2:air-around-us:air}{aire} --- poca
sustancia, poca \omterm{def:g3:measuring-mass:mass}{masa}. La lata es pequeña pero está repleta de
sustancia --- más \omterm{def:g3:measuring-mass:mass}{masa}, así que hunde su platillo. La \omterm{def:g3:measuring-mass:mass}{masa} cuenta
sustancia, no tamaño.
\end{solution}

\begin{solution}{exo:g3:measuring-mass:6}
$\qty{2}{kg} = \qty{2000}{g}$; medio kilogramo son \qty{500}{g};
$\qty{1200}{g}$ es más que $\qty{1}{kg} = \qty{1000}{g}$.
\end{solution}

\begin{solution}{exo:g3:measuring-mass:7}
\qty{100}{g} todavía: partirla cambia la forma, no la sustancia ---
la ley ``misma sustancia, misma \omterm{def:g3:measuring-mass:mass}{masa}'', con todos los trozos juntos en
el platillo.
\end{solution}

\begin{solution}{exo:g3:measuring-mass:8}
La báscula de cocina compara con un muelle apretado que lleva dentro;
la \omterm{def:g3:levers-and-scales:scale}{balanza} de dos platillos compara con pesas a través de brazos
iguales --- la regla de la \omterm{def:g3:levers-and-scales:lever}{palanca} del capítulo de las \omterm{def:g3:levers-and-scales:lever}{palancas}.
\end{solution}

\begin{solution}{exo:g3:measuring-mass:9}
Primera manera: $500 + 200 + 100 = \qty{800}{g}$ de pesas enfrente de
la harina. Segunda manera: poner la pesa de \qty{100}{g} \emph{junto a
la harina} y $500 + 200 + 200 = \qty{900}{g}$ enfrente; unos platillos
horizontales significan entonces que harina $+ \qty{100}{g}$ es igual
a \qty{900}{g}, así que la harina es $900 - 100 = \qty{800}{g}$.
\end{solution}

\begin{solution}{exo:g3:measuring-mass:10}
La ley no ha fallado: no ha desaparecido nada de sustancia --- solo se
ha separado. Los \qty{40}{g} que faltan son la piel. Una pesada más lo
demuestra: la piel sola debería marcar \qty{40}{g}, y
$160 + 40 = \qty{200}{g}$.
\end{solution}

\begin{solution}{exo:g3:measuring-mass:11}
Tienen la \emph{misma} \omterm{def:g3:measuring-mass:mass}{masa} --- un \omterm{def:g3:measuring-mass:units}{kilogramo} cada uno. Las plumas
ocupan muchísimo más sitio. La gente responde mal porque las manos y
los ojos juzgan que el hierro es ``sustancia pesada''; pero la
pregunta ya había fijado la \omterm{def:g3:measuring-mass:mass}{masa} en un \omterm{def:g3:measuring-mass:units}{kilogramo} para cada uno ---
lo único que cambia es el sitio que ocupan.
\end{solution}
